% from ega3-1-fr.tex
\oldpage[0\textsubscript{III}]{63}
ou lorsqu'il existe un entier $m$ tel que tout $\mathcal O_X$-Module
admette une résolution flasque de longueur $\leq m$
(\hyperref[0.11.4.4-fr]{11.4.4}).
\end{env}

\begin{env}[12.4.2]
\label{0.12.4.2-fr}
Nous désignerons de même par $\mathbf H^\bullet(X,\mathcal K^\bullet)$
l'hypercohomologie du foncteur $\Gamma$ par rapport à un complexe
$\mathcal K^\bullet$ de $\mathcal O_X$-Modules~; les
$\mathbf H^p(X,\mathcal K^\bullet)$ sont donc des
$\Gamma(X,\mathcal O_X)$-modules. On peut d'ailleurs considérer
$\mathbf H^\bullet(X,\mathcal K^\bullet)$ comme un cas particulier de
$\mathcal H^\bullet(f,\mathcal K^\bullet)$, où $f$ est un morphisme de
$(X,\mathcal O_X)$ sur un espace annelé réduit à un point muni de l'anneau
$\Gamma(X,\mathcal O_X)$.

Pour tout ouvert $V$ de $X$, nous écrirons
$\mathbf H^\bullet(V,\mathcal K^\bullet)$ au lieu de
$\mathbf H^\bullet(V,\mathcal K^\bullet|V)$.
\end{env}

\begin{proposition}[12.4.3]
\label{0.12.4.3-fr}
\emph{Pour tout entier $p\in\mathbb Z$, le $\mathcal O_Y$-Module
$\mathcal H^p(f,\mathcal K^\bullet)$ est canoniquement isomorphe au
faisceau associé au préfaisceau
$U\mapsto\mathbf H^p(f^{-1}(U),\mathcal K^\bullet)$ sur $Y$.}
\end{proposition}

\begin{proof}
En effet, avec les notations de
(\hyperref[0.12.4.1-fr]{12.4.1}), le faisceau de cohomologie
$\mathcal H^p(f_*(\mathcal L^{\bullet\bullet}))$ est associé au préfaisceau
\[
U\longmapsto H^p(\Gamma(U,f_*(\mathcal L^{\bullet\bullet})))
=H^p(\Gamma(f^{-1}(U),\mathcal L^{\bullet\bullet})).
\]
Mais il est clair que
$\mathcal L^{\bullet\bullet}|f^{-1}(U)$ est une résolution injective de
Cartan-Eilenberg de $\mathcal K^\bullet|f^{-1}(U)$ (T, 3.1.3), donc
\[
H^p(\Gamma(f^{-1}(U),\mathcal L^{\bullet\bullet}))
=\mathbf H^p(f^{-1}(U),\mathcal K^\bullet)
\]
par définition.
\end{proof}

\begin{proposition}[12.4.4]
\label{0.12.4.4-fr}
\emph{L'hypercohomologie $\mathcal H^\bullet(f,\mathcal K^\bullet)$ est
un foncteur cohomologique en $\mathcal K^\bullet$ dans chacun des cas
suivants~:}
\begin{enumerate}
\item[\emph{a)}]
\emph{$\mathcal K^\bullet$ varie dans la catégorie des complexes limités
inférieurement.}
\item[\emph{b)}]
\emph{Il existe un entier $m$ tel que tout $\mathcal O_X$-Module admette
une résolution flasque de longueur $\leq m$.}
\item[\emph{c)}]
\emph{$X$ est un espace noethérien.}
\end{enumerate}
\end{proposition}

\begin{proof}
Les cas a) et b) sont des cas particuliers de
(\hyperref[0.11.5.4-fr]{11.5.4}). D'autre part, le cas c) se déduit de
(\hyperref[0.11.5.2-fr]{11.5.2}), car on sait que dans ce cas, le foncteur
$f_*$ permute aux limites inductives (G, II, 3.10.1).
\end{proof}

\begin{env}[12.4.5]
\label{0.12.4.5-fr}
Considérons maintenant un recouvrement ouvert
$\mathfrak U=(U_\alpha)$ de $X$, et pour tout complexe de
\emph{préfaisceaux} $\mathcal K^\bullet=(\mathcal K^j)$ sur $X$, le
bicomplexe $\check C^\bullet(\mathfrak U,\mathcal K^\bullet)$, dont le
composant d'indices $(i,j)$ est
$\check C^i(\mathfrak U,\mathcal K^j)$, groupe des $i$-cochaînes alternées
du nerf de $\mathfrak U$ à valeurs dans $\mathcal K^j$ (G, II, 5.1). Nous
dirons que la cohomologie de ce bicomplexe est
l'\emph{hypercohomologie du recouvrement $\mathfrak U$ à coefficients dans
$\mathcal K^\bullet$}, et nous la noterons
\[
\mathbf H^\bullet(\mathfrak U,\mathcal K^\bullet)
=H^\bullet(\check C^\bullet(\mathfrak U,\mathcal K^\bullet)).
\]
La suite spectrale de Leray d'un recouvrement (T, 3.8.1 et G, II, 5.9.1)
se généralise comme suit à l'hypercohomologie~:
\end{env}

\begin{proposition}[12.4.6]
\label{0.12.4.6-fr}
\emph{Soit $\mathcal K^\bullet=(\mathcal K^i)$ un complexe de
$\mathcal O_X$-Modules. Il existe un foncteur spectral régulier en
$\mathcal K^\bullet$ ayant pour aboutissement l'hypercohomologie
$\mathbf H^\bullet(X,\mathcal K^\bullet)$, et dont le terme $E_2$ est donné
par}
\begin{equation}
\label{0.12.4.6.1-fr}
\tag{12.4.6.1}
E_2^{pq}=\check H^p(\mathfrak U,h^q(\mathcal K^\bullet)),
\end{equation}
\emph{où $h^q(\mathcal K^\bullet)$ désigne le complexe de préfaisceaux
$V\mapsto H^q(V,\mathcal K^\bullet)$ sur $X$. La suite spectrale précédente
est birégulière si $\mathcal K^\bullet$ est limité inférieurement.}
\end{proposition}

\begin{proof}
Considérons une résolution injective de Cartan-Eilenberg
$\mathcal L^{\bullet\bullet}$ de $\mathcal K^\bullet$, et le tricomplexe
$\check C^\bullet(\mathfrak U,\mathcal L^{\bullet\bullet})
=(\check C^i(\mathfrak U,\mathcal L^{jk}))$~; considérons d'abord ce
tricomplexe comme un bicomplexe pour les degrés $i$ et $j+k$. Comme $i$ ne
prend que des valeurs $\geq0$, la seconde suite spectrale de ce bicomplexe
est régulière (\hyperref[0.11.3.3-fr]{11.3.3}) et dégénérée, car on a
$\check H^q(\mathfrak U,\mathcal L^{jk})=0$ pour tout $q>0$, les
$\mathcal O_X$-Modules $\mathcal L^{jk}$ étant des faisceaux
